\documentclass[11pt]{article}
\usepackage[a4paper,margin=24mm]{geometry}
\usepackage[T1]{fontenc}
\usepackage{lmodern}
\usepackage{microtype}
\usepackage{graphicx}
\usepackage{xcolor}
\usepackage{caption}
\usepackage{enumitem}
\usepackage{float}
\usepackage{fancyhdr}
\usepackage[hidelinks]{hyperref}

\definecolor{Accent}{HTML}{315B7D}
\definecolor{Muted}{HTML}{5B6570}

\captionsetup{font=small,labelfont={bf,color=Accent}}
\setlist[itemize]{leftmargin=1.3em,itemsep=2pt,topsep=3pt}
\newcommand{\Faq}{\textit{F.~aquilonia}}
\newcommand{\Fpo}{\textit{F.~polyctena}}
\newcommand{\pt}[1]{\medskip\noindent\textbf{\color{Accent}#1}\par\nobreak\smallskip}

\pagestyle{fancy}\fancyhf{}
\lhead{\small\color{Muted}The three polyctena blocks}
\rhead{\small\color{Muted}talk notes}
\cfoot{\small\thepage}
\renewcommand{\headrulewidth}{0.4pt}
\setlength{\parskip}{0.4em}\setlength{\parindent}{0pt}

\title{\vspace{-3.2em}\textbf{\color{Accent}The three ``unusual'' polyctena blocks}\\[0.2em]
\large What we can --- and cannot --- say about them, and the Sielva story}
\author{\normalsize Poikela et al. \; -- \; talk-preparation summary}
\date{\normalsize\today}

\begin{document}
\maketitle
\vspace{-1.8em}

\noindent\textit{A short, honest walk through one thread: three large genomic
blocks that behave unlike anything else in the hybrids, why they are tempting to
read as co-adapted (epistasis / clustered incompatibilities), and exactly how far
the data let us take that idea. Colour code: purple = \Faq{}, yellow = \Fpo{}.}

%% ------------------------------------------------------------------
\pt{The observation}
Across the diagnostic loci, sorting toward \Fpo{} is dominated by \textbf{three
large, low-recombination blocks} --- on chromosomes 5, 25 and 26. They are
unusual on three counts at once: they are big, they sit in low recombination,
and they are the \emph{only} strongly \emph{repeatable} \Fpo{} signal --- every
one of the 20 hybrid populations has near-fixed \Fpo{} there. That combination
is exactly what a set of co-adapted / epistatically-required variants would look
like, so it is worth asking how well the data support that reading.

\begin{figure}[H]\centering
\includegraphics[width=0.72\textwidth]{../Figures/di25_popfix_snp.png}
\caption{Per-population near-fixation: 20 population rings, each cell coloured by
that population's own near-fixation direction. The three blocks (Chr 5, 25, 26)
are \textbf{solid yellow across every ring} --- fixed \Fpo{} in all populations ---
whereas \Faq{} fixation is patchier and more distributed.}
\end{figure}

%% ------------------------------------------------------------------
\pt{Can we test co-adaptation directly? (LD between the blocks)}
The clean test would be linkage disequilibrium \emph{between} the blocks: if the
correct alleles are required together, their rare \Faq{} variants should co-occur,
and mismatched combinations should be missing. On the surface that is exactly what
we see --- the all-\Faq{} combination is hugely over-represented and every mixed
class is absent. \textbf{But it is population structure, not individual-level
co-selection.} Only three populations carry any \Faq{} here, each with its own
fixed pattern (\r{A}land at Chr26, Nyrh\"ispera75 at Chr5, Sielva at all three),
so the between-block correlation (Chr5--Chr25 $r=0.83$) collapses to $-0.03$ once
those two populations are removed. The ``missing'' mixed genotypes carry no weight
either: under independence their expected counts are below one anyway.

\begin{figure}[H]\centering
\includegraphics[width=\textwidth]{../Figures/di25_region_dmi.png}
\caption{\textbf{(a)} Every hybrid carrying \Faq{} at these blocks belongs to one
of three populations, each with a fixed regional pattern. \textbf{(b)} The three-block
genotype combinations: the all-\Faq{} class is the five Sielva samples; mixed classes
are absent but expected at $<1$ under independence. The apparent signal is which
population an individual is from.}
\end{figure}

\noindent This is the crux, and it matches the intuition we started with: the blocks
are near-fixed, and LD needs variation --- there simply is not enough segregating
\Faq{} here, and what little there is tracks populations. So co-adaptation cannot be
\emph{tested} in this panel, only motivated.

%% ------------------------------------------------------------------
\pt{Do the blocks segregate within populations?}
A softer prediction of the same idea: if mismatched recombinants are less viable,
each population should be \emph{uniform} at these blocks rather than a segregating
mix. That holds --- within-population variation is essentially zero everywhere
(only two single-block exceptions), so no recombinant genotypes are present. The
interesting outlier is \textbf{Sielva}: the most F1-like population (81\% heterozygous
genome-wide, roughly twice any other), it uniquely retains the ancestral
\emph{heterozygous} state at all three blocks, where every more-sorted population has
fixed \Fpo{}. (Caveat: the five Sielva samples are one colony --- effectively a single
genotype.)

\begin{figure}[H]\centering
\includegraphics[width=\textwidth]{../Figures/di25_region_segregation.png}
\caption{Block ancestry vs each population's genome-wide heterozygosity (F1-ness).
Error bars $\approx 0$: populations are uniform at the blocks (no segregation).
Sielva sits alone at far right --- F1-like and heterozygous at all three blocks.}
\end{figure}

%% ------------------------------------------------------------------
\pt{Is Sielva special --- and has it begun to prune?}
Two things about Sielva turned out to matter. First, it is \emph{not} unusually
``opposite'' anywhere: it is heterozygous at $\sim$84\% of the genome (it is an F1),
and it is \textbf{never} homozygous for the opposite allele against a consensus. So
its state at the three blocks is simply its default --- the blocks are special because
of the \emph{other} 19 populations, not because of Sielva.

Second, and more interesting: even this young, F1-like colony has \textbf{already begun
to prune heterozygosity at sorted loci, directionally}. As loci become more sorted in
the other populations, Sielva's heterozygosity falls from $\sim$83\% to $\sim$52\%, and
the homozygous fraction is almost entirely \emph{toward} the same allele the others fix
(121:1). This survives a recombination control (the effect of sorting is $\sim$10$\times$
stronger than recombination and highly significant). The lone exception is the three
huge blocks themselves --- held as intact F1 heterozygous blocks, because pruning needs
recombination to expose the disfavoured allele, and those blocks resist it longest.

\begin{figure}[H]\centering
\includegraphics[width=\textwidth]{../Figures/di25_sielva_pruning.png}
\caption{Sielva's genotype composition vs how sorted a locus is in the other 19
populations. Heterozygosity (grey) shrinks as loci sort; the homozygous fraction is
almost entirely \emph{toward} the fixed allele (blue), not against (red). Directional
early pruning, beyond what recombination explains.}
\end{figure}

%% ------------------------------------------------------------------
\pt{Take-homes}
\begin{itemize}
  \item Three large low-recombination blocks (Chr 5, 25, 26) are the only
        \textbf{repeatably \Fpo{}} signal --- fixed in all 20 populations. A natural
        candidate for co-adapted / clustered incompatibilities.
  \item That hypothesis \textbf{cannot be tested} here: the blocks are near-fixed, the
        residual \Faq{} variation is confined to three populations, so between-block LD
        is inseparable from population structure and the ``missing'' genotypes carry no
        statistical weight.
  \item The blocks \textbf{do not segregate within populations}, and the one F1-like
        population (Sielva) still holds them as intact heterozygous blocks --- consistent
        with the picture, though Sielva is a single colony.
  \item Even young Sielva shows \textbf{directional early pruning} at ordinary sorted loci
        (losing heterozygosity toward the fixed allele), while shielding the three big
        blocks --- exactly where recombination-sheltered co-adapted variants would persist
        longest.
  \item Net: a coherent and suggestive story, but one the current admixture panel can
        \textbf{motivate, not prove}. Testing it would need populations still segregating
        \Faq{}/\Fpo{} at these blocks, or a functional / mapping approach.
\end{itemize}

\end{document}
